\documentclass{beamer}

\usepackage{aula}

\title{Modelos de linguagem locais}
\date{\today}

\begin{document}
    
    \input{capa}

% Slide 1
\begin{frame}{O que é Ollama?}
    \begin{itemize}
        \item \textbf{Ollama} é uma plataforma para executar LLMs localmente.
        \item Simplifica download, versionamento e execução de modelos.
        \item Fornece API HTTP compatível com aplicações externas.
        \item Foco em privacidade, controle e baixa latência.
    \end{itemize}
\end{frame}

% Slide 2
\begin{frame}{Por que executar LLMs localmente?}
    \begin{itemize}
        \item Privacidade de dados sensíveis.
        \item Redução de custos com APIs externas.
        \item Baixa latência.
        \item Customização e fine-tuning local.
        \item Uso offline.
    \end{itemize}
\end{frame}

% Slide 3
\begin{frame}{Arquitetura Simplificada}
    \begin{itemize}
        \item Cliente (CLI ou aplicação)
        \item Servidor Ollama (API REST)
        \item Modelo baixado (ex: llama3, mistral)
        \item Execução via CPU ou GPU
    \end{itemize}
\end{frame}

% Slide 4
\begin{frame}{Pré-requisitos}
    \begin{itemize}
        \item Docker instalado
        \item Mínimo recomendado:
        \begin{itemize}
            \item 16GB RAM
            \item CPU moderna ou GPU compatível
            \item Espaço em disco para modelos (4GB+)
        \end{itemize}
    \end{itemize}
\end{frame}

% Slide 5
\begin{frame}{Instalação com Docker}

\end{frame}

% Slide 6
\begin{frame}{Baixando um Modelo}
Após iniciar o container:

\begin{verbatim}
docker exec -it ollama ollama pull llama3
\end{verbatim}

Ou via API:
\begin{verbatim}
curl http://localhost:11434/api/pull \
-d '{"name": "llama3"}'
\end{verbatim}
\textbf{1. Baixar a imagem oficial:}
\begin{verbatim}
    docker pull ollama/ollama
\end{verbatim}

\textbf{2. Executar o container:}
\begin{verbatim}
    docker run -d \
    -v ollama:/root/.ollama \
    -p 11434:11434 \
    --name ollama \
    ollama/ollama
    \end{verbatim}\end{frame}

% Slide 7
\begin{frame}{Executando o Modelo}
Via CLI:

\begin{verbatim}
docker exec -it ollama ollama run llama3
\end{verbatim}

Via API:

\begin{verbatim}
curl http://localhost:11434/api/generate \
-d '{
  "model": "llama3",
  "prompt": "Explique LLMs"
}'
\end{verbatim}
\end{frame}

% Slide 8
\begin{frame}{Integração com Aplicações}
    \begin{itemize}
        \item API REST disponível em:
        \begin{center}
        http://localhost:11434
        \end{center}
        \item Pode ser integrada com:
        \begin{itemize}
            \item Python
            \item Node.js
            \item Sistemas corporativos
        \end{itemize}
    \end{itemize}
\end{frame}

% Slide 9
\begin{frame}{Gerenciamento de Modelos}
    \begin{itemize}
        \item Listar modelos:
\begin{verbatim}
ollama list
\end{verbatim}
        \item Remover modelo:
\begin{verbatim}
ollama rm llama3
\end{verbatim}
        \item Versionamento automático
    \end{itemize}
\end{frame}

% Slide 10
\begin{frame}{Boas Práticas}
    \begin{itemize}
        \item Monitorar uso de memória e CPU/GPU
        \item Usar volumes Docker persistentes
        \item Avaliar quantização de modelos
        \item Restringir acesso externo à API
        \item Testar modelos menores antes de escalar
    \end{itemize}
\end{frame}

% Slide Extra
\begin{frame}{Controle do Ollama com Open WebUI}
    
    \textbf{Arquitetura:}
    \begin{itemize}
        \item Navegador → Open WebUI
        \item Open WebUI → API HTTP do Ollama
        \item Ollama → Modelo local (ex: llama3)
    \end{itemize}
    
    \vspace{0.3cm}
    
    \textbf{Executando Open WebUI com Docker:}
    
    \begin{verbatim}
        docker run -d \
        -p 3000:8080 \
        -e OLLAMA_BASE_URL=http://host.docker.internal:11434 \
        --name openwebui \
        ghcr.io/open-webui/open-webui:main
    \end{verbatim}
    
    \vspace{0.3cm}
    
    Acesse:
    \begin{center}
        http://localhost:3000
    \end{center}
    
\end{frame}

\end{document}